% ==========================================
% LatexRefinement Step Output: step4-02-17-25-tuesday-all-offset-final.tex
% Model: gemini-3-flash-preview
% Temperature: 0.33
% TopP: 0.95
% TopK: 40
% MaxOutputTokens: 65535
% ThinkingBudget: 32765
% ThinkingLevel: HIGH
% Prompt Tokens: 36'376
% Candidates Tokens: 13'067 (inkl. Thinking Tokens)
% Cached Tokens: 0
% Processed on: 2026-07-15 09:47:47
% ==========================================

% Primary Language: German

\lecturechapter[1]{Dienstag}{17. Feb}{17. Februar 2026}{Einführung, Organisation und Syntax der Prädikatenlogik}


\begin{nice-box}[Syllabus \& Agenda]
In der heutigen Vorlesung beschäftigen wir uns mit den folgenden Themen:
\begin{itemize}
    \item \textbf{Einführung und Organisation:} Vorstellung des Teams, der Übungsgruppen, des Moodle-Forums und der Leistungsbewertung.
    \item \textbf{Syntax der Prädikatenlogik erster Stufe:}
    \begin{itemize}
        \item Das formale Alphabet (logische und nicht-logische Symbole, Signaturen).
        \item Terme (Präfix- vs. Infix-Notation).
        \item Wohlgebildete Formeln und deren induktiver Aufbau.
    \end{itemize}
    \item \textbf{Freie und gebundene Variablen:} Rekursive Definition und Unterscheidung.
    \item \textbf{Substitution:} Zulässigkeit von Ersetzungen und Variablenkollisionen.
    \item \textbf{Logische Axiome:} Einführung in die aussagenlogischen, Quantoren- und Gleichheitsaxiome.
\end{itemize}
\end{nice-box}

\section{Einführung und Administration}

\begin{spoken-clean}[00:00:00 - 00:01:45]
Herzlich willkommen zur Vorlesung Grundstrukturen. Mein Name ist Christian Urech. Ich arbeite hier als Senior Scientist mit Fokus Education. Nebenbei forsche ich noch in algebraischer Geometrie und geometrischer Gruppentheorie. Und ich freue mich sehr auf dieses Semester mit Ihnen.

Okay, beginnen wir mit dem Anfang --- ich sage das gerne zu Beginn der Vorlesung, dass wir uns mal überlegen: Was sind überhaupt die Ziele? Das erste Ziel erwähnen wir in der Regel nur in der ersten Stunde, aber es ist wichtig, dass wir es immer im Hinterkopf behalten: Das erste Ziel ist auf jeden Fall, gesund und glücklich zu bleiben. 

Viel Mathematik zu lernen ist natürlich sehr wichtig, aber das Ganze hat keinen Sinn, wenn Sie dabei nicht gesund bleiben. Denken Sie also immer daran, egal was geschieht: Das wichtigste Ziel ist, dass Sie in jeglicher Hinsicht gesund bleiben. Und ich möchte Sie gerne dazu ermuntern, zu sich selbst Sorge zu tragen und auch gegenseitig Sorge zu tragen.
\end{spoken-clean}

\subsection{Ziele der Vorlesung}
\begin{math-stroke}
\begin{enumerate}[label=\textbf{(\arabic*)}]
    \item \textbf{Gesund und glücklich bleiben.}
    \item \textbf{Viel Mathematik lernen.}
    \item \textbf{Eine gute Zeit verbringen.}
\end{enumerate}
\end{math-stroke}

\begin{spoken-clean}[00:01:45 - 00:03:15]
Dann eben: Sehr weit oben auf der Liste steht auch noch, viel Mathematik zu lernen. Das ist das Ziel, das wir meistens vor Augen haben werden. Und dann an dritter Stelle --- auch nicht unwichtig --- steht noch, dass wir eine gute Zeit verbringen. Auch das darf man nicht vergessen. 

Ich denke, Sie alle sind hier an die ETH gekommen mit einer gewissen Freude an der Mathematik, mit Ambitionen und mit Hoffnungen. Und ich möchte Sie gerne ermuntern, diese Freude an der Mathematik und diese Hoffnungen und Ambitionen nicht zu vergessen, auch wenn so die große Walze des ganzen Materials über Sie hinwegrollt. Machen Sie vielleicht hin und wieder einmal einen Waldspaziergang und überlegen Sie: Weshalb mache ich das eigentlich? 

Es ist ja eigentlich eine schöne Situation: Sie dürfen vom Morgen bis zum Abend einfach genau das tun, was Ihnen am meisten oder zumindest viel Spaß macht --- nämlich das Fach, das Sie sich ausgesucht haben. Und auch wenn es hart ist oder dann auch sehr schnell zu viel wird, sollte man das nicht vergessen.

Vielleicht noch ein Hinweis: Die ETH bietet auch Ressourcen an, falls Sie eben beim ersten Punkt manchmal Probleme haben. Zögern Sie also nicht, diese Ressourcen auch in Anspruch zu nehmen. Mental-psychische Gesundheit ist ein wichtiges Thema an Hochschulen --- vernachlässigen Sie es nicht.
\end{spoken-clean}

\begin{meta-note}[Projizierter Inhalt: Ressourcen zur psychischen Gesundheit]
Der Dozent zeigt einen Link auf die Beratungsstelle der ETH Zürich für Studierende: \url{https://ethz.ch/studierende/de/beratung/studium-mentale-gesundheit.html}.
\end{meta-note}

\subsection{Organisation der Vorlesung}

\begin{spoken-clean}[00:03:15 - 00:05:10]
Okay, dann noch ein bisschen zur Organisation. Die Vorlesungen finden jeweils am Dienstag von zwei bis vier statt, im G3 --- oder nicht, wir sind im G3, nicht im G5; Entschuldigung, im G3. Aber auf jeden Fall in einem dieser schönen Hörsäle mit Fenstern. Da bin ich froh, dass wir im Frühling nicht in den Keller gehen müssen.

Ah, warten Sie, das ist Quatsch, da steht noch Mittwoch. Okay, das kann man vergessen, es ist am Dienstagnachmittag im G3. Die verlässlichen Informationen --- ich habe das von der letzten Vorlesung übernommen und dachte, ich hätte es korrigiert, aber nicht --- und Dokumente finden Sie auf der Moodle-Seite der Vorlesung. Falls Sie also noch nicht auf Moodle sind, gehen Sie dorthin und suchen Sie alle Sachen. Dort finden Sie auch die Übungsblätter, dort können Sie die Übungen abgeben, und Sie finden alle Informationen, die Sie wahrscheinlich für diese Vorlesung brauchen, und mehr.

Es gibt eine schriftliche Prüfung in der Prüfungssession. Die Note ist zu 100 Prozent die Note, die Sie an der Prüfung haben. Sie dürfen also machen, was Sie wollen über das Semester; Sie müssen einfach die Prüfung schreiben --- oder dürfen die Prüfung schreiben, wenn Sie wollen ---, und die Note, die Sie dort haben, ist dann Ihre Note.
\end{spoken-clean}

\begin{math-stroke}[Organisation der Vorlesung]
\begin{itemize}
    \item \textbf{Vorlesungen:} Dienstag 14:15 -- 16:00, HG G3.
    \item \textbf{Moodle:} Alle Informationen und Dokumente zur Vorlesung auf Moodle.
    \item \textbf{Prüfung:} Schriftliche Prüfung in der Prüfungssession ($100$\% der Endnote).
    \item \textbf{Skript:} Skript von Lorenz Halbeisen.
\end{itemize}
\end{math-stroke}

\begin{spoken-clean}[00:05:10 - 00:07:30]
Wir folgen dem Skript von Professor Lorenz Halbeisen. Er ist der Logiker im Haus, Titularprofessor hier für Logik. Er hat wesentlich dazu beigetragen, diese Vorlesung Grundstrukturen zu konzipieren und aufzustellen, vor etwa fünf Jahren oder so. Und wir folgen diesem Skript. Wenn es also inhaltliche Beschwerden gibt, dann dürfen Sie sich natürlich an mich wenden --- ich übernehme die Verantwortung dafür. Das Skript ist auf der Moodle-Seite; es ist ein gutes Skript und von der Fachperson im Haus gemacht.

Es gibt noch weitere Bücher; es gibt quasi das Buch zum Skript. Lorenz Halbeisen hat zusammen mit Regula Krapf dieses umfassendere Einführungsbuch in die Logik geschrieben: \qt{Gödel's Theorems and Zermelo's Axioms}. Falls Sie noch etwas mehr Tiefe wollen, als das Skript bietet, dann können Sie in diesem Buch zum Beispiel nachlesen. 

Aber es gibt selbstverständlich noch ganz viele andere Lehrbücher über Logik. Ich habe noch einen Link auf der Moodle-Seite eingefügt: Es gibt \qt{Logic Matters}, das ist ein Blog von einem Logikprofessor --- ich glaube aus Cambridge oder Oxford ---, und dort gibt es sehr, sehr viele Referenzen. Man findet auch viele Blogs und alles Mögliche im Internet. Also: Logik ist gut.
\end{spoken-clean}

\begin{math-stroke}[Literatur]
\begin{itemize}
    \item \textbf{Hauptquelle:} Skript von Lorenz Halbeisen.
    \item \textbf{Ergänzende Literatur:} L. Halbeisen, R. Krapf: \qt{Gödel's Theorems and Zermelo's Axioms} (Birkhäuser-Verlag 2020).
    \item \textbf{Online-Ressourcen:} Blog \qt{Logic Matters} (\url{https://www.logicmatters.net/}).
\end{itemize}
\end{math-stroke}

\subsection{Organisation der Übungen}

\begin{spoken-clean}[00:07:30 - 00:09:50]
Gut, dann zur Organisation der Übungen. Die Übungen finden jeweils am Mittwoch statt, natürlich in ganz unterschiedlichen Räumen. Das können Sie nachschauen; es kommt darauf an, für welche Übungsgruppe Sie sich eingeschrieben haben. Der Übungskoordinator ist Konstantin Andritsch. Bei Fragen zu den Übungen können Sie ihm direkt eine E-Mail schreiben und sich an ihn wenden.

Es gibt sieben Übungsgruppen. Wahrscheinlich haben Sie sich bereits für eine dieser Gruppen eingeschrieben. Wenn Sie das noch nicht getan haben, tun Sie das bitte. Und dann gehen Sie bitte auch in die Übungsgruppe, für die Sie sich eingeschrieben haben. Die Übungen, die Sie abgeben, gehen automatisch an die Assistierenden Ihrer Übungsgruppe. Das heißt, da kann man keine Wechsel machen. Eine der Gruppen ist auf Englisch; wenn Sie lieber Englisch haben, dann können Sie dorthin gehen, oder auch wenn Sie Englisch lernen wollen --- wobei die meisten von Ihnen wahrscheinlich sowieso genügend gut Englisch sprechen, sodass es keine Rolle spielt. Eine dieser Gruppen ist noch in Form einer Fokusgruppe, das kennen Sie wahrscheinlich bereits aus dem ersten Semester.

Es gibt jede Woche eine Übungsserie. Die kommt jeweils am Dienstag ungefähr heraus, vielleicht manchmal schon am Montag, im schlimmsten Fall am Dienstagabend oder so. Dann haben Sie bis zum Montagmorgen der nächsten Woche Zeit, diese zu lösen. Spätestens dann müssen Sie sie abgeben; Sie dürfen natürlich gerne auch schon früher abgeben. Die Assistierenden geben Ihnen dann am Mittwoch in der Stunde Feedback zu den Übungen. Es gibt jetzt bereits eine erste Serie. Da können Sie morgen schon ein bisschen Fragen stellen und diese mit Ihren Assistierenden vorbesprechen, weil morgen bereits Übungsstunden sind, und dann nächste Woche abgeben.
\end{spoken-clean}

\begin{math-stroke}[Organisation der Übungen]
\begin{itemize}
    \item \textbf{Termin:} Mittwoch 14:15 -- 16:00.
    \item \textbf{Koordination:} Konstantin Andritsch.
    \item \textbf{Gruppen:} $7$ Übungsgruppen (eine auf Englisch, eine Fokusgruppe).
    \item \textbf{Ablauf:} Wöchentliche Serien (Ausgabe Dienstag, Abgabe Montag 08:00 via Moodle).
\end{itemize}
\end{math-stroke}

\begin{spoken-clean}[00:09:50 - 00:12:10]
Einfach nochmals zur Erinnerung --- das haben Ihnen wahrscheinlich schon sehr viele Professoren gesagt: Übungen sind wirklich sehr wichtig. Sie sollten die Übungen machen und abgeben, und wir gehen davon aus, dass Sie das tun. Das ist ein allgemein sehr wichtiges Problem: Oft schaut man sich den Stoff an und denkt: \qt{Ah ja, das ist alles klar, das ist okay, ich habe das verstanden.} Vielleicht gibt es solche, die gehen sogar die Übungen durch und denken: \qt{Ah, ich weiß, wie man das löst, ich weiß, wie man das löst.} Aber vielleicht wissen Sie nicht, wie man es aufschreibt. Dann gehen Sie an die Prüfung, schreiben irgendetwas hin, und natürlich wird bei der Prüfung bewertet, was Sie aufschreiben, und nicht, was Sie gemeint haben. Dann bekommen Sie plötzlich ganz viele Punkte abgezogen, weil Sie halt nicht sehr sinnvolle Sachen hingeschrieben haben, obwohl Sie das Richtige gemeint haben.

Es ist also sehr wichtig, dass Sie auch lernen, mathematisch aufzuschreiben, und das tun Sie, indem Sie Übungen abgeben. Übungen abzugeben ist ein Riesenservice, den Sie hier erhalten: Die ETH bezahlt viele Assistierende, die stundenlang Ihre Übungen korrigieren und Ihnen persönliches Feedback geben. Ich möchte Sie wirklich ermuntern, davon Gebrauch zu machen. 

Mathematik --- ja, das haben Ihnen wahrscheinlich schon viele gesagt --- Mathematik ist nichts zum Nur-Zuschauen. Lernen kann man es nur, indem man es selbst macht. Das ist wie Fußballspielen oder Geigespielen: Man kann noch so viele Fußballspiele schauen --- wenn man dann selbst auf dem Platz steht, ist man wahrscheinlich noch nicht so gut im Spielen.
\end{spoken-clean}

\begin{didactic-insight}[Die Bedeutung des mathematischen Aufschreibens]
Der Dozent betont einen zentralen Aspekt des Mathematikstudiums: Die Diskrepanz zwischen dem intuitiven Verständnis einer Lösung und der Fähigkeit, diese formal korrekt und präzise niederzuschreiben. Die Übungen dienen nicht nur der inhaltlichen Vertiefung, sondern sind das primäre Training für die mathematische Kommunikation und Beweisführung, welche in der Prüfung bewertet wird.
\end{didactic-insight}

\begin{spoken-clean}[00:12:10 - 00:14:15]
Insbesondere in dieser Vorlesung ist das wichtig. Wenn Sie schauen: Das ist ja nur eine zweistündige Vorlesung, es ist keine so riesengroße Vorlesung. Aber es gibt fünf Credits. Für eine größere Vorlesung wie Lineare Algebra bekommen Sie sieben Credits. Das heißt, wenn man es auf den Schlüssel herunterbricht, erhalten Sie für die Übungen dieser Vorlesung genauso viele Credits wie für die Übungen der Linearen Algebra. 

Wir erwarten auch, dass Sie für die Übungen dieser Vorlesung etwa gleich viel Zeit aufwenden wie für die Übungen der Linearen Algebra. Das heißt: Diese Vorlesung hier ist viel übungsbasierter als andere Vorlesungen. Es wird trotzdem lange Übungsblätter geben, und Sie sollten bitte auch viel Zeit verwenden, um diese zu bearbeiten. Davon gehen wir aus, und insbesondere gehen wir davon aus, dass Sie viel Zeit mit den Übungen verbracht haben, wenn wir die Prüfung vorbereiten. Aber es ist auch schön, Übungen zu machen; man versteht dann die Sachen endlich richtig.

Okay, dann noch zu den Software-Tools. Wir werden immer wieder mit Clicker-Fragen arbeiten --- nicht heute, aber in späteren Wochen. Wenn wir eine Frage stellen, können Sie mit der EduApp abstimmen, ob es richtig oder falsch ist, welche Auswahl zutrifft, und so können wir es ein bisschen interaktiv machen. Es ist immer ein bisschen die Frage: Wie gestaltet man so eine große Vorlesung interaktiv und persönlich, sodass jeder mitmachen und sich irgendwie beteiligen kann? Mit verschiedenen Software-Tools versucht man, eine gewisse Interaktivität zu kreieren, und eines davon sind eben diese Clicker-Fragen.
\end{spoken-clean}

\subsection{Software-Tools}
\begin{math-stroke}
\begin{itemize}
    \item \textbf{EduApp:} Für interaktive Clicker-Fragen während der Vorlesung.
    \item \textbf{Kursforum:} Auf Moodle zur Diskussion inhaltlicher Fragen.
\end{itemize}
\end{math-stroke}

\begin{spoken-clean}[00:14:15 - 00:16:30]
Das Zweite ist --- und dazu möchte ich Sie auch sehr stark ermuntern --- dieses Kursforum auf Moodle. Wir haben ein Forum auf Moodle, wo Sie Fragen zur Vorlesung stellen und auch Fragen beantworten können. Ich möchte Sie wirklich ermuntern, das zu verwenden. Ich kann Ihnen zeigen, wie das geht; das ist wirklich so wie --- vielleicht kennen Sie es --- \qt{Stack Exchange} oder so, wo man Fragen stellen kann. Das hier ist eines speziell nur für diese Vorlesung. Das heißt, Sie können hier mit Ihren Kommilitoninnen über die Vorlesung diskutieren. 

Gehen wir hier auf die Moodle-Seite: Da haben wir alle Informationen, das Skript und ein paar Links. Und dann gehen wir hier zum Forum zu Grundstrukturen. Das Ganze ist anonym. Das heißt --- bitte verhalten Sie sich zivilisiert, aber wir sind ja alle erwachsene Menschen ---, Sie können wirklich ohne Angst Fragen stellen. Sie müssen keine Sorge haben, dass irgendjemand denkt: \qt{Oh nein, das ist eine blöde Frage oder eine blöde Antwort.} Sie können da wirklich frei von der Leber weg Fragen stellen. 

Man kann hier sagen: \qt{Add a new discussion topic}. Ich weiß auch nicht, dann können wir zum Beispiel fragen: \qt{Ganze Zahlen}. Und dann die Frage: \qt{Gibt es die ganzen Zahlen noch, wenn alles Leben ausgestorben ist?} Eine interessante Frage. Dann kann man \qt{Post to forum} machen. Jetzt können Sie dorthin gehen, und jemand anderes kann diese Frage beantworten. Es gibt viele Antworten, man kann das diskutieren und Rückfragen stellen.
\end{spoken-clean}

\begin{meta-note}[Demonstration des Moodle-Forums]
Der Dozent zeigt live im Browser die Moodle-Seite der Vorlesung und demonstriert das Erstellen eines neuen Diskussionsthemas im Forum. Er betont die Anonymität und die Möglichkeit zur gegenseitigen Unterstützung.
\end{meta-note}

\begin{spoken-clean}[00:16:30 - 00:18:50]
Man kann auch sagen: Das ist eine gute Frage, und man kann sie hochvoten. Seien Sie vorsichtig mit Downvoten, das ist nicht sehr nett --- also lieber nur hochvoten, wenn Sie finden: \qt{Ah ja, diese Frage hatte ich auch}, oder wenn es eine gute Antwort ist. Die Assistierenden werden auch immer ein Auge auf das Forum haben, um zu schauen, dass keine falschen Antworten überhandnehmen. Auch wir werden, wenn eine Antwort richtig ist, sie als solche markieren. Das heißt, Sie haben dann eine quasi garantiert korrekte Antwort, der Sie vertrauen können.

Ich glaube, das ist aus verschiedenen Gründen sehr nützlich. Einerseits bleibt man manchmal einfach bei einer Frage stecken, und dann ist es besser, man fragt direkt jemanden; es gibt viele hier, die diese Frage beantworten können. Andererseits ist das Beantworten selbst ein sehr wichtiger Prozess. Es ist fast noch besser als Übungen zu lösen, Fragen in Foren zu beantworten. Übungen sind oft ein bisschen künstlich, oder? Man hat eine Frage, die jemand gestellt hat, weiß aber, dass die Person, die die Antwort lesen wird, es vielleicht besser verstanden hat als man selbst. 

Hier ist es anders: Man muss die Antwort so formulieren, dass die Person, die die Frage gestellt hat, sie versteht --- und dass trotzdem alles korrekt ist. Das ist eine sehr gute Übung. Es ist auch eine gute Übung, Fragen zu stellen. Und es ist immer netter und besser, Fragen an andere Menschen zu stellen anstatt nur an ChatGPT --- obwohl Sie dort auch oft gute Antworten kriegen, aber ein Forum ist trotzdem die bessere Variante. Ich glaube, es ist wichtig, dass man einfach mal ins Laufen kommt. Springen Sie über Ihren Schatten, stellen Sie einfach einmal eine Frage oder beantworten Sie eine, und dann gibt das mit der Zeit hoffentlich einen regen Betrieb. Gut, so viel zum Forum --- also wirklich als Motivation.
\end{spoken-clean}

\begin{ai-global-state-checkpoint-invisible-content}
% timestamp: 00:18:50
% topic: Administrative Einführung und Forum-Demo
% board_state: none (slides only)
% next_goal: Inhaltsübersicht der Vorlesung
% open_loops: none
\end{ai-global-state-checkpoint-invisible-content}

\subsection{Inhalt der Vorlesung}

\begin{spoken-clean}[00:18:50 - 00:21:15]
Gut, jetzt zum Inhalt der Vorlesung. Es ist eine --- ich würde sagen --- gewissermaßen spezielle Vorlesung, vielleicht nicht ganz Standard. In Analysis, Linearer Algebra oder Gruppentheorie wird an fast allen Unis weltweit sehr ähnliches Material unterrichtet. Eine Vorlesung wie \qt{Grundstrukturen} gibt es nicht an allen Unis. Sie ist auch hier an der ETH relativ neu; ich glaube, es gibt sie seit etwa fünf Jahren. Die Idee ist, ein paar wirklich grundlegende Sachen zu besprechen, für die man in anderen Vorlesungen keine oder nur wenig Zeit hat.

Im ersten Kapitel beginnen wir ganz am Anfang mit der Logik. Da geht es wirklich darum, die Mathematik von Grund auf aufzubauen. Das ist gar nicht so einfach. Sie denken jetzt vielleicht: \qt{Okay, das haben wir doch bereits in Linearer Algebra und Analysis gemacht.} Dort hat man vielleicht schon mehr von Grund auf angefangen, als Sie es aus der Mittelschule für möglich gehalten hätten, aber man steigt trotzdem noch recht weit oben ein. Hier beginnen wir nun mit der Prädikatenlogik erster Stufe. Da fangen wir wirklich ganz am Anfang an, das Ganze logisch aufzubauen.
\end{spoken-clean}

\begin{math-stroke}[Inhalt der Vorlesung]
\begin{itemize}
    \item \textbf{Prädikatenlogik erster Stufe:} Logischer Aufbau der Mathematik von Grund auf.
    \item \textbf{Zermelo-Fraenkel Mengenlehre (ZF):} Das theoretische Fundament der modernen Mathematik.
    \item \textbf{Konstruktion der reellen Zahlen:} Präzise Definition von $\mathbb{R}$ (Dedekindsche Schnitte).
    \item \textbf{Auswahlaxiom (AC):} Ein spezielles und fundamentales Axiom der Mengenlehre.
    \item \textbf{Kardinalzahlen:} Vergleich von Unendlichkeiten.
    \item \textbf{Weitere Themen:} Graphentheorie, elementare Zahlentheorie, etc.
\end{itemize}
\end{math-stroke}

\begin{spoken-clean}[00:21:15 - 00:23:45]
Man muss sich etwas daran gewöhnen; es ist nicht ganz einfach, zu entscheiden: Wo beginnt man jetzt wirklich? Aber Sie werden sehen. Es ist wichtig, dass Sie einen Einblick erhalten: Wie geht das überhaupt? Was sind diese logischen Aussagen? Was sind Beweise? Was sind Axiome, und wie verwendet man sie? Wir werden dann die Zermelo-Fraenkel-Axiome anschauen. Das sind die üblichen Axiome, auf denen \qt{theoretisch} zumindest ein Großteil der modernen Mathematik aufbaut. Ich sage \qt{theoretisch}, weil sehr wenige Mathematiker ihre Beweise wirklich bis auf die Axiome zurückführen --- man steigt meist viel weiter oben ein.

Es ist trotzdem eine Mischung; es ist kein reiner Logik-Kurs. Dafür ist viel zu wenig Zeit. Es ist nur der erste Teil der Vorlesung, und es ist nur eine zweistündige Vorlesung. Wenn man das Ganze sauber, gründlich und in allen Details machen möchte, bräuchte man ein ganzes Semester lang eine vierstündige Vorlesung nurfür diese Themen --- das ist eigentlich das, was das Buch hier macht. Da es im zweiten Semester keine volle Logik-Vorlesung braucht, werden wir auf gewisse Details nicht zu stark insistieren, damit wir weiterkommen und Sie einfach einen Eindruck erhalten, wie das funktioniert.
\end{spoken-clean}

\begin{spoken-clean}[00:23:45 - 00:26:00]
Dann werden wir noch ein paar Sachen machen, die wichtig sind und für die man in anderen Vorlesungen keine Zeit hatte: zum Beispiel die Konstruktion der reellen Zahlen. Das hatten Sie in der Analysis nur am Rande gestreift. Was sind die reellen Zahlen überhaupt, und wie kann man sie konstruieren? Dann werden wir das Auswahlaxiom anschauen --- ein spezielles Axiom der Zermelo-Fraenkel-Mengenlehre. Dann schauen wir uns Kardinalzahlen an. 

In einem zweiten Teil geht es darum, dass wir wirklich konkrete Mathematik machen. Da geht es um ein bisschen Graphentheorie und elementare Zahlentheorie --- einfach Themen, in die Sie einen Einblick erhalten sollen, für die in anderen Vorlesungen nicht so viel Zeit ist. Auch dabei geht es vor allem wieder darum, dass Sie lernen: Wie funktioniert mathematisches Begründen? Wie führt man Beweise? Es ist also auch da wieder wichtig, dass Sie die Übungen machen.

Das ist der Inhalt der Vorlesung. Vielleicht noch eine letzte Bemerkung dazu, was \emph{nicht} Inhalt der Vorlesung ist: Wir beginnen zwar ganz grundlegend mit der Logik, aber wir beginnen nicht auf \qt{Stufe Null} --- das wäre dort, wo das Ganze in die Philosophie übergeht. Da müsste man sich überlegen: Was ist Mathematik überhaupt?
\end{spoken-clean}

\begin{math-stroke}[Nicht Teil der Vorlesung]
\begin{itemize}
    \item \textbf{Philosophie der Mathematik:} Die Frage nach dem Wesen der Mathematik (Stufe Null) wird nicht vertieft behandelt.
\end{itemize}
\end{math-stroke}

\begin{spoken-clean}[00:26:00 - 00:28:30]
Wir beginnen jetzt mit der Prädikatenlogik erster Stufe. Das sieht so aus: Wir bauen ein System auf, eine Sprache, in der wir eigentlich einfach Zeichen aneinanderreihen. Dann stellen wir Regeln auf, wie wir diese Zeichen aneinanderreihen können, was dann einen Beweis ergibt, und das schauen wir uns im Vergleich zur Mathematik an. Das ist kein philosophisches Thema; das macht man einfach so. 

Aber man kann natürlich sofort philosophisch fragen: Ist Mathematik jetzt einfach nur diese Aneinanderreihung von Zeichen? Oder ist diese Aneinanderreihung von Zeichen einfach ein sehr nützliches Werkzeug, um Mathematik zu betreiben? Existieren Zahlen überhaupt? Und wenn ja, in welcher Form? Existieren vielleicht nur endliche Zahlen wirklich und alles andere ist reiner Formalismus? Oder ist Mathematik nur etwas, das das menschliche Denken erschafft, und außerhalb des menschlichen Denkens existiert sie gar nicht? Was genau ist es? Das sind viele spannende philosophische Fragen, über die man schon seit Jahrtausenden diskutiert.

Diese Logik, wie wir sie verwenden, ist stark vom Formalismus inspiriert --- das war eben Hilbert und so weiter. Aber man kann heute genauso gut sagen, dass das eher einfach nur ein Werkzeug ist und die Mathematik nochmals etwas anderes. Das sind schwierige philosophische Fragen und leider nicht Teil dieser Vorlesung --- aber es ist ja auch kein Philosophiestudium, sondern ein Mathematikstudium.
\end{spoken-clean}

\begin{spoken-clean}[00:28:30 - 00:30:30]
Vielleicht gibt es doch manche unter Ihnen, die das interessiert. Da möchte ich Sie ermuntern, sich einmal ein paar Bücher anzuschauen und sich Gedanken zu machen --- das wird kein Teil der Prüfung sein. Gerade wenn man Logik betreibt, ist das ein sehr guter Moment, um sich die ganz grundlegenden Fragen zu stellen: Was machen wir da überhaupt? Was ist überhaupt Mathematik?

Es gibt sehr viel Literatur zur Philosophie der Mathematik. Da ist zum einen dieses Buch auf Deutsch: \qt{Einführung in die Philosophie der Mathematik} von Jörg Neunhäuserer. Es ist sehr kurz und verständlich aus der Sicht eines Mathematikers geschrieben. Ich glaube, seine eigenen Kommentare sind philosophisch nicht extrem tiefschürfend, aber seine Erklärungen sind sehr hilfreich, und man erhält in relativ kurzer Zeit einen Einblick in die verschiedenen Standpunkte. Dann gibt es ein anderes, das sehr empfohlen wird --- ich habe es noch nicht ganz gelesen ---, nämlich von Shapiro: \qt{Thinking about Mathematics}. Das ist auch sehr schön und lesbar geschrieben. Das wären Quellen.

Und dann noch eine nicht obligatorische Aufgabe für dieses Semester: Setzen Sie sich doch einfach einmal in ein Café oder eine Bar mit Ihren Kommilitoninnen oder anderen Leuten und diskutieren Sie darüber: Was ist eigentlich Mathematik? Aber wie gesagt: Das ist nicht Teil dieser Vorlesung, sondern nur eine Nebenbemerkung.
\end{spoken-clean}

\subsection{Literatur}
\begin{math-stroke}
\begin{itemize}
    \item \textbf{J. Neunhäuserer:} \qt{Einführung in die Philosophie der Mathematik}.
    \item \textbf{S. Shapiro:} \qt{Thinking about Mathematics}.
\end{itemize}
\end{math-stroke}

\begin{spoken-clean}[00:30:30 - 00:31:29]
Gut, so weit zur Einführung und zur Information. Gibt es gerade noch Fragen, die Ihnen unter den Nägeln brennen? Ansonsten können Sie immer noch E-Mails schreiben oder im Forum Fragen stellen --- das ist noch besser. Ja?

\inlinemetanote{Ein Student stellt eine Frage zur Nachfolgevorlesung.}

Die Nachfolgevorlesung von Grundstrukturen? Das hängt ein bisschen von den Kapiteln ab. Diese Sachen hier sind wirklich die Basics, die man eigentlich kennen sollte. Für Algebra braucht man überall das Auswahlaxiom; da müssen Sie wissen, was das ist, oder die verschiedenen Kardinalitäten kennen. Das sind absolute Grundlagen. 

Was die Logik betrifft, gibt es direkt keine Nachfolgevorlesung, aber Lorenz Halbeisen bietet immer mal wieder eine Logikvorlesung an --- nicht regelmäßig, aber alle paar Jahre --- oder ein Seminar zum Thema Logik. Das wäre die natürliche Fortsetzung, ist aber nicht obligatorisch. Ansonsten sind quasi alle algebraischen und zahlentheoretischen Vorlesungen Nachfolgevorlesungen. Es sind eben Grundlagen, die man eigentlich für alles braucht.

Kein Notenbonusprogramm, nein. Einfach nur das, was ich gesagt habe: die Prüfung. Gut, okay. Sonst steht wie gesagt das Forum offen. Dann beginnen wir jetzt.
\end{spoken-clean}

\begin{meta-note}[Tafelübergang]
Der Dozent schaltet den Projektor aus, trinkt einen Schluck Wasser und bereitet die Tafel vor, um mit dem ersten inhaltlichen Kapitel der Vorlesung zu beginnen.
\end{meta-note}

\section{Syntax}

\begin{spoken-clean}[00:31:29 - 00:31:45]
So. Was wir heute machen, erscheint vielleicht teilweise noch etwas seltsam, aber wir beginnen jetzt mit dem Kapitel Null im Skript: Syntax.
\end{spoken-clean}

\begin{math-stroke}
\textbf{0. Syntax}
\end{math-stroke}

\subsection{Das Alphabet}

\begin{spoken-clean}[00:31:45 - 00:32:57]
Hier sind es wirklich nur Zeichen auf der syntaktischen Ebene. Später werden das dann Variablen sein, die beispielsweise für Zahlen stehen, wenn man in der Zahlentheorie arbeitet, oder für Mengen, wenn man in der Mengenlehre arbeitet. Wenn wir in der Gruppentheorie arbeiten, stehen sie für Elemente von Gruppen, und wenn wir Lineare Algebra machen, für Vektoren. Aber egal --- hier sind es einfach nur Zeichen. \inlinemetanote{schreibt an die Tafel}
\end{spoken-clean}

\begin{math-stroke}[Das Alphabet: Variablen]
\begin{enumerate}[label=\textbf{(\alph*)}]
    \item \textbf{Variablen:} z.\,B. $x, y, v_0, v_1, \dots$
\end{enumerate}
\end{math-stroke}

\begin{spoken-clean}[00:32:57 - 00:34:00]
Das Zweite sind logische Operatoren. \inlinemetanote{schreibt an die Tafel} Davon gibt es vier: Es gibt diesen Haken, der heißt \qt{nicht} ($\neg$). Dann gibt es ein Dach nach oben, das ist \qt{und} ($\wedge$). Ein nach oben offener Keil, das ist \qt{oder} ($\vee$). Und dann gibt es noch das Zeichen --- einfach ein Pfeil ---, das heißt \qt{impliziert} ($\rightarrow$). 

Die Namen dieser Zeichen sind natürlich bereits sehr suggestiv, und es ist auch klar, wie wir diese später interpretieren werden, aber auch hier gilt: A priori sind das nur Zeichen.
\end{spoken-clean}

\begin{math-stroke}[Logische Operatoren]
\begin{enumerate}[label=\textbf{(\alph*)}]
    \setcounter{enumi}{1}
    \item \textbf{Logische Operatoren:} 
    \begin{itemize}
        \item $\neg$ (nicht)
        \item $\wedge$ (und)
        \item $\vee$ (oder)
        \item $\rightarrow$ (impliziert)
    \end{itemize}
\end{enumerate}
\end{math-stroke}

\begin{spoken-clean}[00:34:00 - 00:35:25]
Gut, dann haben wir \textbf{(c)}, das wären die logischen Quantoren. \inlinemetanote{schreibt an die Tafel} Da gibt es zwei: \qt{es existiert} --- das ist der Existenzquantor ($\exists$) --- und \qt{für alle} --- der Allquantor ($\forall$). 

Unter \textbf{(d)} gibt es noch ein Relationszeichen, und das ist die Gleichheitsrelation, also das Gleichheitszeichen ($=$). Das Zeichen heißt \qt{gleich}. Diese Symbole \textbf{(a)}, \textbf{(b)}, \textbf{(c)} und \textbf{(d)} heißen Logiksymbole. 

Nun gibt es noch nicht-logische Symbole, die wir ebenfalls verwenden werden. Das wäre unter \textbf{(e)}: Konstantensymbole.
\end{spoken-clean}

\begin{math-stroke}[Logische Quantoren und Gleichheit]
\begin{enumerate}[label=\textbf{(\alph*)}]
    \setcounter{enumi}{2}
    \item \textbf{Logische Quantoren:}
    \begin{itemize}
        \item $\exists$ (Existenzquantor)
        \item $\forall$ (Allquantor)
    \end{itemize}
    \item \textbf{Gleichheitsrelation:} $=$
\end{enumerate}
\end{math-stroke}

\begin{spoken-clean}[00:35:25 - 00:37:35]
Das ist, wie wir sehen werden, theoriespezifisch. Die Symbole \textbf{(a)} bis \textbf{(d)} haben wir immer, egal in welcher Theorie wir arbeiten. Dazu kommen dann noch theoriespezifische Symbole. Konstantensymbole zum Beispiel --- machen wir dazu auch Beispiele, wobei das auch wieder einfach nur Symbole sind: In der Zahlentheorie haben wir das Symbol $0$. Wenn Sie Mengenlehre betreiben, haben wir die leere Menge. Das ist ein Symbol für die leere Menge, ein bestimmtes Objekt --- einfach ein Konstantensymbol. 

Gut, dann gibt es noch Funktionssymbole. Auch das sind wieder Symbole, die wir uns gewissermaßen als Funktionen denken, aber auch das hängt von der Theorie ab. In der Zahlentheorie gibt es zum Beispiel das Funktionssymbol $+$, oder in der Analysis den $\sin$.
\end{spoken-clean}

\begin{math-stroke}[Nicht-logische Symbole: Konstanten und Funktionen]
\begin{enumerate}[label=\textbf{(\alph*)}]
    \setcounter{enumi}{4}
    \item \textbf{Konstantensymbole:} z.\,B. $0$ in der Zahlentheorie, oder $\emptyset$ in der Mengenlehre.
    \item \textbf{Funktionssymbole:} z.\,B. $+$ in der Zahlentheorie oder $\sin$ in der Analysis.
\end{enumerate}
\end{math-stroke}

\begin{spoken-clean}[00:37:35 - 00:40:30]
Was bei Funktionssymbolen noch wichtig ist: Sie haben eine Stelligkeit. Ein Funktionssymbol $F$ hat eine \newterm{Stelligkeit} --- das ist einfach die Anzahl der Argumente, die dieses Funktionssymbol hat. Was wäre zum Beispiel die Stelligkeit von $+$? Ja?
\end{spoken-clean}

\begin{student-interaction}[Student Answer]
Zwei.
\end{student-interaction}

\begin{spoken-clean}[continued]
Zwei, genau. Man kann hier und dort etwas einsetzen. Der $\sin$ hat dagegen nur ein Argument. 

Dann gibt es noch Relationssymbole, das ist Punkt \textbf{(g)}. Auch da wieder ein Beispiel: In der Zahlentheorie haben wir \qt{kleiner gleich} oder \qt{strikt kleiner}. Man hat zwei Zahlen, und diese erfüllen die Relation oder eben nicht. Das sind einfach Relationssymbole. In der Mengenlehre haben wir \qt{ist enthalten}, also: $x$ ist Element von $y$. Das ist auch ein Relationssymbol. 

Auch Relationssymbole haben eine Stelligkeit. Was sind das für Stelligkeiten bei \qt{kleiner} und \qt{Element}? Ja?
\end{spoken-clean}

\begin{student-interaction}[Student Answer]
Zwei.
\end{student-interaction}

\begin{spoken-clean}[continued]
Zwei, genau.
\end{spoken-clean}

\begin{math-stroke}[Relationssymbole und Stelligkeit]
\begin{enumerate}[label=\textbf{(\alph*)}]
    \setcounter{enumi}{6}
    \item \textbf{Relationssymbole:} z.\,B. $<$ in der Zahlentheorie, oder $\in$ in der Mengenlehre.
\end{enumerate}

\begin{explanation-of-steps}[Stelligkeit von Symbolen]
Jedes Funktions- oder Relationssymbol besitzt eine feste \newterm{Stelligkeit} (Arität), welche die Anzahl der Argumente festlegt:
\begin{itemize}
    \item \textbf{Funktionen:} $+$ hat Stelligkeit $2$, $\sin$ hat Stelligkeit $1$.
    \item \textbf{Relationen:} $<$ und $\in$ haben Stelligkeit $2$.
\end{itemize}
\end{explanation-of-steps}
\end{math-stroke}

\begin{spoken-clean}[00:40:30 - 00:41:45]
Die Symbole \textbf{(a)} bis \textbf{(d)} heißen logische Symbole, \textbf{(e)} bis \textbf{(g)} nicht-logische Symbole. Wenn wir so eine Menge von Symbolen für eine Theorie haben, dann nennen wir die Menge aller nicht-logischen Symbole die \newterm{Signatur}. 

Gut, das sind erst einmal die Zeichen, die wir haben. Als Nächstes wollen wir festlegen, was ein Term ist. Zeichen kann man ja einfach beliebig aneinanderreihen --- das ergibt dann Reihen von Zeichen. Aber jetzt definieren wir, was ein Term ist.
\end{spoken-clean}

\begin{math-stroke}[Vollständige Übersicht: Das formale Alphabet]
\textbf{Das formale Alphabet der Prädikatenlogik erster Stufe:}
\begin{enumerate}[label=\textbf{(\alph*)}]
    \item \textbf{Variablen:} z.\,B. $x, y, v_0, v_1, \dots$
    \item \textbf{Logische Operatoren:} $\neg$ (nicht), $\wedge$ (und), $\vee$ (oder), $\rightarrow$ (impliziert)
    \item \textbf{Logische Quantoren:} $\exists$ (Existenzquantor), $\forall$ (Allquantor)
    \item \textbf{Gleichheitsrelation:} $=$
    \item \textbf{Konstantensymbole:} z.\,B. $0, \emptyset$ (theorieabhängig).
    \item \textbf{Funktionssymbole:} z.\,B. $+$, $\sin$ (mit zugehöriger Stelligkeit).
    \item \textbf{Relationssymbole:} z.\,B. $<, \in$ (mit zugehöriger Stelligkeit).
\end{enumerate}

Die Symbole \textbf{(a)}--\textbf{(d)} heißen \newterm{logische Symbole}.
Die Symbole \textbf{(e)}--\textbf{(g)} heißen \newterm{nicht-logische Symbole}.

\begin{definition}[Signatur]
\label[definition]{def:signatur}
Die \newterm{Signatur} $\mathcal{L}$ ist die Menge der nicht-logischen Symbole.
\end{definition}
\end{math-stroke}

\subsection{Terme}

\begin{spoken-clean}[00:41:45 - 00:43:25]
Dafür sagen wir jetzt: Sei $\mathcal{L}$ eine Signatur, also eine Menge von nicht-logischen Symbolen. Ein $\mathcal{L}$-Term ist eine Zeichenkette, die durch endlich viele Anwendungen der folgenden Regeln entstanden ist. 

Was ist also ein Term? Die erste Regel ist \textbf{(T0)} --- \qt{T} wie Term: Jede Variable ist ein $\mathcal{L}$-Term. Oft sagen wir auch einfach nur \qt{Term}, wenn aus dem Kontext klar ist, mit welcher Signatur wir arbeiten. Wir wissen also: Jede Variable an sich ist schon ein $\mathcal{L}$-Term. 

Dann \textbf{(T1)}: Sinnvollerweise legt man fest, dass auch jede Konstante ein $\mathcal{L}$-Term ist.
\end{spoken-clean}

\begin{math-stroke}[Induktive Definition von Termen (Teil 1)]
Sei $\mathcal{L}$ eine Signatur. Ein \newterm{$\mathcal{L}$-Term} ist eine Zeichenkette, die durch endlich viele Anwendungen der folgenden Regeln entstanden ist:
\begin{enumerate}[label=\textbf{(T\arabic*)}, start=0]
    \item Jede Variable ist ein $\mathcal{L}$-Term.
    \item Jede Konstante $c \in \mathcal{L}$ ist ein $\mathcal{L}$-Term.
\end{enumerate}
\end{math-stroke}

\begin{spoken-clean}[00:43:25 - 00:45:05]
Dann haben wir noch \textbf{(T2)}: Wenn wir schon Terme haben\dots \inlinemetanote{schreibt an die Tafel} Seien $\tau_1$ bis $\tau_n$ $\mathcal{L}$-Terme. Wenn $F$ ein $n$-stelliges Funktionssymbol in $\mathcal{L}$ ist, dann ist $F \tau_1 \dots \tau_n$ auch wieder ein $\mathcal{L}$-Term. 

So, wie wir es jetzt machen, schreiben wir einfach $F$ und hängen hinten die Argumente dran ($F \tau_1 \dots \tau_n$). Das braucht keine Klammern und so weiter --- ich werde das nach der Pause noch etwas genauer ausführen. Aber so, wie wir es meistens schreiben, ist es eigentlich $F(\tau_1, \dots, \tau_n)$. 

Man kann also einfach Variablen und Konstanten nehmen, darauf Funktionen anwenden und dann natürlich auch wieder Funktionen von Funktionen von Variablen und Konstanten bilden und so weiter. Das ist dann jeweils wieder ein $\mathcal{L}$-Term. Man kann endlich viele davon nehmen und sie wieder in eine neue Funktion einsetzen --- oder in dieselbe --- und erhält so weitere $\mathcal{L}$-Terme.
\end{spoken-clean}

\begin{math-stroke}[Induktive Definition von Termen (Teil 2)]
\begin{enumerate}[label=\textbf{(T\arabic*)}]
    \setcounter{enumi}{1}
    \item Sind $\tau_1, \dots, \tau_n$ $\mathcal{L}$-Terme und ist $F \in \mathcal{L}$ ein $n$-stelliges Funktionssymbol, dann ist die Zeichenkette $F \tau_1 \dots \tau_n$ ebenfalls ein $\mathcal{L}$-Term.
\end{enumerate}

\begin{explanation-of-steps}[Struktur von Termen]
Terme repräsentieren die \qt{Objekte} einer mathematischen Theorie. Die Regel \textbf{(T2)} erlaubt die rekursive Konstruktion beliebig komplexer Ausdrücke. In der formalen Syntax (Präfix-Notation) werden keine Klammern benötigt, da die Stelligkeit von $F$ eindeutig festlegt, wie viele der folgenden Terme als Argumente zu interpretieren sind.
\end{explanation-of-steps}
\end{math-stroke}

\begin{spoken-clean}[00:45:05 - 00:46:40]
Gut, wir machen nach der Pause weiter; wir haben jetzt noch bis Viertel nach Pause. \inlinemetanote{Pause. Der Dozent wischt die Tafel.}
\end{spoken-clean}

\begin{ai-global-state-checkpoint-invisible-content}
% timestamp: 00:46:40
% topic: Definition von Alphabet und Termen
% board_state: def:signatur, def:l-term
% next_goal: Beispiele für Terme, Präfix- vs. Infix-Notation, Definition von Formeln
% open_loops: none
\end{ai-global-state-checkpoint-invisible-content}

\begin{spoken-clean}[00:46:40 - 00:47:30]
\inlinemetanote{Ein Student stellt eine Frage.}
\end{spoken-clean}

\begin{student-interaction}[Studentenfrage]
Entschuldigung, bei den Symbolen, kann man \textbf{(d)} und \textbf{(g)} zusammennehmen?
\end{student-interaction}

\begin{spoken-clean}[continued]
Genau, die Gleichheitsrelation und Relationssymbole\dots Die Gleichheitsrelation ist eine Art von Relation, ja. Aber sie ist einfach so wichtig und kommt in jeder Theorie vor, deswegen führen wir sie separat auf. Die Symbole unter \textbf{(d)} haben wir in jeder Theorie, während \textbf{(g)} theorieabhängig ist. Das Gleichheitszeichen brauchen wir überall, aber es ist im Grunde auch eine Art zweistelliges Relationszeichen. Gut.
\end{spoken-clean}

\begin{spoken-clean}[00:47:30 - 00:49:25]
\inlinemetanote{Der Dozent putzt die Tafel weiter.} Also, putzen wir --- etwas weniger gründlich, aber\dots Wir haben jetzt \textbf{(T0)} und \textbf{(T1)} definiert. Jetzt müssen wir noch \textbf{(T2)} festhalten. Machen wir das noch, damit wir wissen, was ein Term ist.
\end{spoken-clean}

\begin{math-stroke}[Vollständige Definition: $\mathcal{L}$-Term]
\begin{definition}[$\mathcal{L}$-Term]\label[definition]{def:l-term-vollstaendig}
Sei $\mathcal{L}$ eine Signatur. Die Menge der \newterm{$\mathcal{L}$-Terme} ist induktiv definiert durch:
\begin{enumerate}[label=\textbf{(T\arabic*)}, start=0]
    \item Jede Variable $\nu$ ist ein $\mathcal{L}$-Term.
    \item Jedes Konstantensymbol $c \in \mathcal{L}$ ist ein $\mathcal{L}$-Term.
    \item Sind $\tau_1, \dots, \tau_n$ $\mathcal{L}$-Terme und ist $F \in \mathcal{L}$ ein $n$-stelliges Funktionssymbol, dann ist $F \tau_1 \dots \tau_n$ ein $\mathcal{L}$-Term.
\end{enumerate}
\end{definition}
\end{math-stroke}

\begin{spoken-clean}[00:49:25 - 00:51:29]
Gut, wir machen jetzt weiter. Wir haben gesehen --- wir haben definiert ---, was Terme sind. Schauen wir uns doch noch Beispiele an. Aber wie gesagt: Die Beispiele sind eigentlich gar nicht so illustrativ, weil es wirklich so banal ist, wie es klingt. 

Wenn $x$ und $y$ Variablen sind und $F$ und $G$ Funktionssymbole mit Stelligkeit 1 respektive 2, dann sind zum Beispiel Terme: $Fx$. Da $x$ eine Variable ist, ist es ein Term; $F$ ist eine Funktion, also ist $Fx$ wieder ein Term. Und $Gxy$ ist ebenfalls ein Term. Oder $GFxy$ --- das schreibt man dann so. Es sieht grauenhaft aus, das so zu schreiben. Was man hier eigentlich sich vorstellt, ist $F(x)$, dort $G(x, y)$ und hier dann $G(F(x), y)$.
\end{spoken-clean}


\begin{math-stroke}[Beispiele für Terme]
Seien $x, y$ Variablen, $F$ ein $1$-stelliges und $G$ ein $2$-stelliges Funktionssymbol.
\begin{examples}[Term-Konstruktionen]
\begin{itemize}
    \item $F x$ (entspricht $F(x)$)
    \item $G x y$ (entspricht $G(x, y)$)
    \item $G F x y$ (entspricht $G(F(x), y)$)
\end{itemize}
\end{examples}
\end{math-stroke}

\begin{spoken-clean}[00:51:29 - 00:55:59]
Das hier ist die sogenannte \newterm{polnische Notation} oder \newterm{Präfix-Notation}, und das andere ist die sogenannte \newterm{Infix-Notation}. Zur polnischen Notation: Bis etwa zum Anfang des letzten Jahrhunderts gab es eine große Blütezeit der Mathematik und insbesondere der Logik in Polen. Die dortigen Logiker kamen mit dieser Notation auf, daher der Name. 

Der Vorteil der Präfix-Notation --- und das ist der logische Vorteil --- ist natürlich: Man braucht tatsächlich keine Klammern. Man kann alle Formeln ohne Klammern schreiben, und es ist immer eindeutig klar, wie sie gemeint sind. Um die ganzen Definitionen sauber aufzubauen, verwenden wir daher die Präfix-Notation. Wenn wir aber praktisch damit arbeiten wollen, ist der Vorteil der Infix-Notation schlicht die bessere Lesbarkeit.
\end{spoken-clean}

\begin{math-stroke}[Notation: Präfix vs. Infix]
\begin{center}
\begin{tabular}{l | l}
    \toprule
    \textbf{Präfix-Notation (Polnische Notation)} & \textbf{Infix-Notation (Standard)} \\ 
    \midrule
    $F x$     & $F(x)$ \\
    $G x y$    & $G(x, y)$ \\
    $G F x y$   & $G(F(x), y)$ \\
    \bottomrule
\end{tabular}
\end{center}

\begin{explanation-of-steps}[Vergleich der Notationen]
\begin{itemize}
    \item \textbf{Vorteil Präfix:} Benötigt keine Klammern, da die Struktur durch die feste Stelligkeit der Symbole eindeutig bestimmt ist.
    \item \textbf{Vorteil Infix:} Wesentlich bessere Lesbarkeit für Menschen. Wir werden im Verlauf der Vorlesung meist die Infix-Notation verwenden.
\end{itemize}
\end{explanation-of-steps}
\end{math-stroke}

\begin{spoken-clean}[00:55:59 - 00:57:29]
Ich möchte jetzt gar nicht ganz genau definieren, was Infix- und was Präfix-Notation ist. Man kann das zwar sauber definieren, aber dann verbrät man wieder eine Stunde. Ich glaube, es ist klar, was gemeint ist: Bei Präfix können Sie einfach alles aneinanderschreiben, und es ergibt eindeutig einen Sinn. Bei Infix muss man dann eben Klammern und Kommas setzen. Das stört uns ja nicht wirklich --- wir müssten dann eben in unserem Alphabet auch noch Kommas und Klammern aufführen und festlegen, welche Operatoren stärker binden. Das kehren wir jetzt einfach ein bisschen unter den Teppich. Wir sagen einfach, wir wissen, was Infix-Notation ist, und verwenden sie oft. 

Noch deutlicher wird es bei Relationssymbolen. Machen wir dazu auch Beispiele. Wenn wir zum Beispiel $x + y$ schreiben wollen: $+$ ist ein zweistelliges Funktionssymbol; in polnischer Notation wäre das $+xy$. Das ist Präfix, und $x + y$ ist Infix. Dasselbe gilt für Relationssymbole und auch für logische Symbole. In polnischer Notation schreibt man zum Beispiel $=xy$, in Infix-Notation $x = y$. Oder $\wedge xy$ wird zu $x \wedge y$.
\end{spoken-clean}

\begin{math-stroke}[Beispiele: Operatoren und Relationen]
\begin{itemize}
    \item \textbf{Addition:} $+ x y$ (Präfix) $\leftrightarrow x + y$ (Infix)
    \item \textbf{Gleichheit:} $= x y$ (Präfix) $\leftrightarrow x = y$ (Infix)
    \item \textbf{Logisches Und:} $\wedge x y$ (Präfix) $\leftrightarrow x \wedge y$ (Infix)
\end{itemize}
\end{math-stroke}

\subsection{Formeln}

\begin{spoken-clean}[00:57:29 - 00:59:59]
Jetzt kommen wir aber zur Formel. Terme haben wir definiert, jetzt folgen die Formeln. Eine wohlgebildete $\mathcal{L}$-Formel ist eine Zeichenkette, die durch endlich viele Anwendungen der folgenden Regeln entstanden ist. 

Wir stellen jetzt wieder --- wie bei den Termen --- Regeln auf, wie man quasi induktiv Formeln bilden kann. Wir haben \textbf{(F0)}: Wenn $\tau_1$ und $\tau_2$ $\mathcal{L}$-Terme sind, dann ist $=\tau_1\tau_2$ --- in Infix-Notation also $\tau_1 = \tau_2$ --- wiederum eine $\mathcal{L}$-Formel. Ein Term gleich ein anderer Term ergibt uns also eine Formel. 

Dann \textbf{(F1)}, das ist ähnlich bei den theoriespezifischen Relationssymbolen: Sind $\tau_1$ bis $\tau_n$ $\mathcal{L}$-Terme und ist $R$ ein $n$-stelliges Relationssymbol in $\mathcal{L}$, dann ist $R\tau_1\dots\tau_n$ ebenfalls eine $\mathcal{L}$-Formel. \inlinemetanote{schreibt an die Tafel} Das ist relativ einfach und noch keine Induktion; sobald wir Terme und diese Relationssymbole haben, ergeben sich daraus Formeln.
\end{spoken-clean}

\begin{math-stroke}[Induktive Definition von Formeln (Teil 1)]
Eine (wohlgebildete) \newterm{$\mathcal{L}$-Formel} ist eine Zeichenkette, die durch endlich viele Anwendungen der folgenden Regeln entstanden ist:
\begin{enumerate}[label=\textbf{(F\arabic*)}, start=0]
    \item Sind $\tau_1, \tau_2$ $\mathcal{L}$-Terme, dann ist $= \tau_1 \tau_2$ (geschrieben als $\tau_1 = \tau_2$) eine $\mathcal{L}$-Formel.
    \item Sind $\tau_1, \dots, \tau_n$ $\mathcal{L}$-Terme und ist $R \in \mathcal{L}$ ein $n$-stelliges Relationssymbol, dann ist $R \tau_1 \dots \tau_n$ eine $\mathcal{L}$-Formel.
\end{enumerate}
\end{math-stroke}

\begin{spoken-clean}[00:59:59 - 01:02:57]
Jetzt kommt \textbf{(F2)}, und damit die Induktion: Falls $\varphi$ eine $\mathcal{L}$-Formel ist, dann ist $\neg\varphi$ auch wieder eine $\mathcal{L}$-Formel. \inlinemetanote{schreibt an die Tafel} Das heißt, wir können bereits bestehende Formeln negieren, um neue zu erhalten.

Und \textbf{(F3)}: Sind $\varphi$ und $\psi$ $\mathcal{L}$-Formeln, so sind $\wedge\varphi\psi$ --- in Infix-Notation $(\varphi \wedge \psi)$ --- sowie $\vee\varphi\psi$ und $\rightarrow\varphi\psi$ ebenfalls $\mathcal{L}$-Formeln. \inlinemetanote{schreibt an die Tafel} Wir können also Formeln mit logischen Operatoren verknüpfen. Auch hier verwenden wir in der formalen Definition die Präfix-Notation, aber in der Praxis schreiben wir meistens die Infix-Notation mit Klammern.

Dann noch \textbf{(F4)}, was die Quantoren betrifft: Falls $\varphi$ eine $\mathcal{L}$-Formel ist und $\nu$ eine beliebige Variable, dann sind $\exists\nu\varphi$ und $\forall\nu\varphi$ auch wieder $\mathcal{L}$-Formeln. \inlinemetanote{schreibt an die Tafel} Damit können wir über Objekte quantifizieren. 

Das sind Formeln. Wir können nun Terme nach diesen Regeln aneinandersetzen und erhalten Formeln. Wenn wir eine Formel haben, können wir die Regeln \textbf{(F2)} bis \textbf{(F4)} immer wieder anwenden und erhalten so immer komplexere Formeln. Vielleicht noch eine Bemerkung: Die Regeln \textbf{(F0)} und \textbf{(F1)} ergeben eine bestimmte Art von Formeln, die wir \newterm{atomar} nennen. \inlinemetanote{schreibt an die Tafel} Das sind sozusagen die Grundbausteine, die noch keine logischen Verknüpfungen oder Quantoren enthalten.
\end{spoken-clean}

\begin{math-stroke}[Vollständige Definition: $\mathcal{L}$-Formel]
\begin{definition}[$\mathcal{L}$-Formel]\label[definition]{def:l-formel-vollstaendig}
Die Menge der \newterm{$\mathcal{L}$-Formeln} ist induktiv definiert durch:
\begin{enumerate}[label=\textbf{(F\arabic*)}, start=0]
    \item Sind $\tau_1, \tau_2$ $\mathcal{L}$-Terme, dann ist $\tau_1 = \tau_2$ eine $\mathcal{L}$-Formel.
    \item Sind $\tau_1, \dots, \tau_n$ $\mathcal{L}$-Terme und $R \in \mathcal{L}$ ein $n$-stelliges Relationssymbol, dann ist $R \tau_1 \dots \tau_n$ eine $\mathcal{L}$-Formel.
    \item Ist $\varphi$ eine $\mathcal{L}$-Formel, dann ist $\neg \varphi$ eine $\mathcal{L}$-Formel.
    \item Sind $\varphi, \psi$ $\mathcal{L}$-Formeln, dann sind $(\varphi \wedge \psi)$, $(\varphi \vee \psi)$ und $(\varphi \rightarrow \psi)$ $\mathcal{L}$-Formeln.
    \item Ist $\varphi$ eine $\mathcal{L}$-Formel und $\nu$ eine Variable, dann sind $\exists \nu \varphi$ und $\forall \nu \varphi$ $\mathcal{L}$-Formeln.
\end{enumerate}
\end{definition}

\begin{remark}[Atomare Formeln]
Formeln, die ausschließlich durch die Regeln \textbf{(F0)} und \textbf{(F1)} gebildet werden, heißen \newterm{atomar}.
\end{remark}
\end{math-stroke}

\begin{ai-global-state-checkpoint-invisible-content}
% timestamp: 01:02:57
% topic: Induktive Definition von Formeln
% board_state: def:l-formel-vollstaendig
% next_goal: Bereich eines Quantors, freie und gebundene Variablen
% open_loops: none
\end{ai-global-state-checkpoint-invisible-content}

\subsection{Freie und gebundene Variablen}

\begin{spoken-clean}[01:02:57 - 01:05:27]
Es gibt noch einen weiteren Begriff: Wenn $\varphi$ eine Formel ist, die durch \textbf{(F4)} konstruiert wurde --- wir haben also eine Formel und wenden \textbf{(F4)} an ---, dann ist sie von der Form $\exists\nu\psi$ oder $\forall\nu\psi$, wobei $\nu$ eine Variable und $\psi$ eine Formel ist. Dann heißt diese Formel $\psi$, die direkt nach dem $\nu$ kommt, der \newterm{Bereich} des Quantors $\exists$ oder $\forall$. \inlinemetanote{schreibt an die Tafel}

Das ist wichtig, um zu verstehen, worauf sich ein Quantor bezieht. Wenn wir zum Beispiel $\forall x (x = y)$ schreiben, dann ist $(x = y)$ der Bereich des Allquantors.
\end{spoken-clean}

\begin{math-stroke}[Bereich eines Quantors]
\begin{definition}[Bereich eines Quantors]\label[definition]{def:bereich-quantor}
Sei $\varphi$ eine $\mathcal{L}$-Formel der Form $\exists \nu \psi$ oder $\forall \nu \psi$. Dann heißt die Teilformel $\psi$ der \newterm{Bereich} (oder Wirkungsbereich) des entsprechenden Quantors.
\end{definition}
\end{math-stroke}

\begin{spoken-clean}[01:05:27 - 01:08:27]
Falls $\nu$ in diesem Bereich vorkommt, so heißt sie \newterm{gebunden} durch den entsprechenden Quantor. Entsprechend heißt eine Variable, die nicht gebunden ist, \newterm{frei}. \inlinemetanote{schreibt an die Tafel}

Man muss hier ein bisschen aufpassen: Eine Variable kann in einer Formel sowohl frei als auch gebunden vorkommen. Zum Beispiel ist in der Formel $(x = y \wedge \forall x (x = z))$ das erste Vorkommen von $x$ frei, während das zweite Vorkommen im Bereich des Allquantors liegt und somit gebunden ist. Wir bezeichnen mit $\operatorname{frei}(\varphi)$ die Menge der freien Variablen in der Formel $\varphi$.
\end{spoken-clean}

\begin{math-stroke}[Variablen in einem Term]
\begin{definition}[Variablen in einem Term]\label[definition]{def:variablen-in-term}
Für einen $\mathcal{L}$-Term $\tau$ ist die Menge $\operatorname{var}(\tau)$ der in $\tau$ vorkommenden Variablen rekursiv wie folgt definiert:
\begin{enumerate}[label=\textbf{(VT\arabic*)}, start=0]
    \item Ist $\tau \equiv \nu$ eine Variable, so ist $\operatorname{var}(\nu) := \{\nu\}$.
    \item Ist $\tau \equiv c$ eine Konstante, so ist $\operatorname{var}(c) := \emptyset$.
    \item Ist $\tau \equiv F \tau_1 \dots \tau_n$, wobei $F$ ein $n$-stelliges Funktionssymbol ist, so ist:
    \[
    \operatorname{var}(F \tau_1 \dots \tau_n) := \operatorname{var}(\tau_1) \cup \dots \cup \operatorname{var}(\tau_n)
    \]
\end{enumerate}
\end{definition}
\end{math-stroke}

\begin{math-stroke}[Freie und gebundene Variablen]
\begin{definition}[Freie und gebundene Variablen]\label[definition]{def:frei-gebunden}
\begin{itemize}
    \item Ein Vorkommen einer Variablen $\nu$ in einer Formel $\varphi$ heißt \newterm{gebunden}, falls es im Bereich eines Quantors $\exists \nu$ oder $\forall \nu$ (\cref{def:bereich-quantor}) liegt.
    \item Ein Vorkommen einer Variablen $\nu$ heißt \newterm{frei}, falls es nicht gebunden ist.
    \item Wir bezeichnen mit $\operatorname{frei}(\varphi)$ die Menge aller Variablen, die in $\varphi$ mindestens ein freies Vorkommen besitzen (bezogen auf die Menge $\operatorname{var}(\tau)$ der Variablen in ihren Termen, vgl. \cref{def:variablen-in-term}). \inlinemetanote{[Anmerkung: Die Menge $\operatorname{frei}(\varphi)$ lässt sich auch präzise rekursiv definieren: $\operatorname{frei}(\tau_1 = \tau_2) := \operatorname{var}(\tau_1) \cup \operatorname{var}(\tau_2)$, $\operatorname{frei}(R \tau_1 \dots \tau_n) := \operatorname{var}(\tau_1) \cup \dots \cup \operatorname{var}(\tau_n)$, $\operatorname{frei}(\neg \varphi) := \operatorname{frei}(\varphi)$, $\operatorname{frei}(\varphi \circ \psi) := \operatorname{frei}(\varphi) \cup \operatorname{frei}(\psi)$ für $\circ \in \{\wedge, \vee, \to\}$, und $\operatorname{frei}(\exists \nu \varphi) := \operatorname{frei}(\forall \nu \varphi) := \operatorname{frei}(\varphi) \setminus \{\nu\}$.]}
\end{itemize}
\end{definition}
\end{math-stroke}

\begin{ai-global-state-checkpoint-invisible-content}
% timestamp: 01:08:27
% topic: Freie und gebundene Variablen
% board_state: def:bereich-quantor, def:frei-gebunden
% next_goal: Definition eines Satzes, Substitution
% open_loops: none
\end{ai-global-state-checkpoint-invisible-content}

\begin{spoken-clean}[01:08:27 - 01:11:27]
Zum Begriff \qt{Menge}: Ich habe das Wort vorhin schon verwendet, und jemand hat in der Pause kurz danach gefragt. Man muss hier ein bisschen aufpassen: Wenn wir hier von Mengen sprechen, meinen wir Mengen im naiven Sinn. Die Frage ist ja immer: Wenn man solche Definitionen aufschreibt, muss man irgendwo anfangen. Man kann nicht alles von Grund auf definieren. Das geht dann in Richtung Metamathematik --- so wie Metaphysik. Es gibt gewisse naive Begriffe, die man einfach voraussetzt. Einer davon ist eben der Mengenbegriff im naiven Sinn --- also einfach Zusammenfassungen von Objekten, wie etwa die Menge aller freien Variablen.

Das sind aber keine Mengen, von denen wir die Potenzmenge bilden, über die wir quantifizieren oder mit denen wir komplizierte Durchschnitte bilden. Es dient einfach nur der Beschreibung. Dasselbe gilt für das, was ich vorhin sagte: Wenn wir etwa von $\tau_1$ bis $\tau_n$ sprechen, wobei $n$ eine ganze Zahl ist. Wir haben die ganzen Zahlen hier gar nicht definiert --- das kann man an dieser Stelle auch gar nicht ---, aber man hat metamathematisch einen naiven Begriff davon, was endliche ganze Zahlen sind. Und das verwenden wir, um das System aufzubauen.
\end{spoken-clean}

\begin{didactic-insight}[Metamathematik und naive Mengenlehre]
Der Dozent reflektiert über das fundamentale Problem der Zirkularität beim Aufbau der Mathematik: Um die formale Mengenlehre (ZFC) zu definieren, benötigt man bereits eine Sprache und grundlegende Konzepte wie \qt{Menge} oder \qt{natürliche Zahl}. Diese werden auf der \newterm{Meta-Ebene} als intuitiv gegeben vorausgesetzt (\newterm{naive Mengenlehre}), um darauf das formale System zu errichten.
\end{didactic-insight}

\begin{spoken-clean}[01:11:27 - 01:13:27]
Gut, das sind also die freien Variablen. Man könnte diese Definitionen von Bereich und freien Variablen noch viel ausführlicher machen, aber es ist, wie gesagt, keine reine Logikvorlesung.

Eine Formel $\varphi$ heißt ein \newterm{Satz}, falls in ihr keine freien Variablen vorkommen. \inlinemetanote{schreibt an die Tafel} Das sind dann Aussagen, die einen festen Wahrheitswert haben können, sobald wir eine Interpretation festgelegt haben.
\end{spoken-clean}

\begin{math-stroke}[Definition eines Satzes]
\begin{definition}[Satz]\label[definition]{def:satz}
Eine $\mathcal{L}$-Formel $\varphi$ heißt ein \newterm{Satz}, falls sie keine freien Variablen enthält, d.\,h. falls $\operatorname{frei}(\varphi) = \emptyset$.
\end{definition}
\end{math-stroke}

\subsection{Substitution}

\begin{spoken-clean}[01:13:27 - 01:16:27]
Sei $\varphi$ eine Formel, $\nu$ eine Variable und $\tau$ ein Term. Was wir nun tun können: Wir können überall dort, wo $\nu$ in $\varphi$ frei vorkommt, die Variable $\nu$ durch den Term $\tau$ ersetzen. \inlinemetanote{schreibt an die Tafel}

Dieser Prozess heißt \newterm{Substitution}. Wir substituieren $\nu$ durch $\tau$ und schreiben dafür $\varphi(\nu/\tau)$. Das ist eine neue Formel, die wir durch diese Ersetzung erhalten.
\end{spoken-clean}

\begin{math-stroke}[Substitution]
\begin{definition}[Substitution]\label[definition]{def:substitution}
Sei $\varphi$ eine $\mathcal{L}$-Formel, $\nu$ eine Variable und $\tau$ ein $\mathcal{L}$-Term. Die Formel $\varphi(\nu/\tau)$ entsteht aus $\varphi$, indem jedes \emph{freie} Vorkommen (\cref{def:frei-gebunden}) von $\nu$ durch den Term $\tau$ ersetzt wird.
\end{definition}
\end{math-stroke}

\begin{ai-global-state-checkpoint-invisible-content}
% timestamp: 01:16:27
% topic: Definition eines Satzes und Substitution
% board_state: def:satz, def:substitution
% next_goal: Zulässige Substitution, Beispiele
% open_loops: none
\end{ai-global-state-checkpoint-invisible-content}

\begin{spoken-clean}[01:16:27 - 01:19:27]
Nun gibt es gewisse Substitutionen, die zulässig sind, und andere, die es nicht sind. Das hat wieder mit dem Begriff der freien und gebundenen Variablen zu tun. Im Term $\tau$ können ja neue Variablen vorkommen oder solche, die bereits in $\varphi$ enthalten sind. Was wir nicht wollen, ist, dass Variablen aus $\tau$ beim Einsetzen in $\varphi$ dort gebunden werden. \inlinemetanote{schreibt an die Tafel}

Wir sagen daher: Eine Substitution ist \newterm{zulässig}, wenn keine Variable in $\tau$ durch Quantoren in $\varphi$ \qt{eingefangen} wird. Das verhindert, dass sich die Bedeutung der Variablen im Term $\tau$ durch die Substitution ungewollt ändert.
\end{spoken-clean}

\begin{math-stroke}[Zulässige Substitution]
\begin{definition}[Zulässige Substitution]\label[definition]{def:zulaessige-substitution}
Eine Substitution $\varphi(\nu/\tau)$ (\cref{def:substitution}) heißt \newterm{zulässig}, falls keine Variable in $\tau$ (\cref{def:variablen-in-term}) an einer Stelle für $\nu$ substituiert wird, die im Bereich eines Quantors dieser Variablen liegt (d.\,h. keine Variable in $\tau$ wird durch die Substitution \qt{eingefangen}, vgl. \cref{def:frei-gebunden}).
\end{definition}
\end{math-stroke}

\begin{spoken-clean}[01:19:27 - 01:22:27]
Schauen wir uns Beispiele an, das macht es vielleicht klarer. \inlinemetanote{zeichnet eine Tabelle an die Tafel} Wir haben hier die Formel $\varphi$, die Variable $\nu$, den Term $\tau$, die substituierte Formel $\varphi(\nu/\tau)$ und die Frage, ob das Ganze zulässig ist.

Nehmen wir zum Beispiel $\varphi \equiv x = y$. Wenn wir $x$ durch eine Konstante $c$ ersetzen, erhalten wir $c = y$. Das ist natürlich zulässig. Wenn wir $x$ durch $y$ ersetzen, erhalten wir $y = y$. Auch das ist zulässig, da es keine Quantoren gibt.

Aber schauen Sie sich das hier an: $\varphi \equiv \forall y (x = y)$. Wenn wir hier $x$ durch $y$ ersetzen, erhalten wir $\forall y (y = y)$. Das ist \emph{nicht} zulässig! Warum? Weil die Variable $y$ im Term $\tau$ plötzlich durch den Allquantor in $\varphi$ gebunden wird. Die ursprüngliche Formel sagte aus, dass $x$ gleich jedem $y$ ist. Nach der Substitution steht dort, dass jedes $y$ gleich sich selbst ist. Das ist eine völlig andere Aussage!
\end{spoken-clean}

\begin{math-stroke}[Beispiele für Substitutionen]
\begin{center}
\begin{tabular}{l | c | c | l | l}
    \toprule
    \textbf{Formel $\varphi$} & \textbf{Var. $\nu$} & \textbf{Term $\tau$} & \textbf{Subst. $\varphi(\nu/\tau)$} & \textbf{Zulässig?} \\ 
    \midrule
    $x = y$ & $x$ & $c$ & $c = y$ & ja \\
    $x = y$ & $x$ & $y$ & $y = y$ & ja \\
    $x = x \wedge \forall x (x = x)$ & $x$ & $c$ & $c = c \wedge \forall x (x = x)$ & ja \\
    $\forall y (x = y)$ & $x$ & $y$ & $\forall y (y = y)$ & \textcolor{BrickRed}{\textbf{nein}} \\
    \bottomrule
\end{tabular}
\end{center}

\begin{explanation-of-steps}[Vermeidung von Variablen-Einfang]
Im letzten Beispiel wird die Variable $y$ des Terms $\tau$ durch den Quantor $\forall y$ \qt{eingefangen}. Durch Umbenennen der gebundenen Variablen (z.\,B. $\forall z (x = z)$) ließe sich dieser Konflikt vermeiden.
\end{explanation-of-steps}
\end{math-stroke}

\begin{ai-global-state-checkpoint-invisible-content}
% timestamp: 01:22:27
% topic: Beispiele für Substitutionen
% board_state: tab:substitution-examples
% next_goal: Logische Axiome (Beamer)
% open_loops: none
\end{ai-global-state-checkpoint-invisible-content}

\subsection{Logische Axiome}

\begin{spoken-clean}[01:22:27 - 01:25:27]
Gut, das Letzte, was ich heute noch kurz ansprechen möchte, sind die logischen Axiome. Das machen wir jetzt per Beamer. \inlinemetanote{schaltet den Projektor ein}

Bisher haben wir nur Zeichen, Terme und Formeln eingeführt, aber das hat a priori noch keinen Sinn. Was wir jetzt brauchen, ist eine Liste von Axiomen --- ein bisschen im Stile Hilberts. Das geht auf David Hilbert zurück, der das große Programm hatte, die Mathematik auf ganz solide Füße stellen zu wollen. Wir haben nun eine Reihe von ausgezeichneten Formeln.

Die erste Serie von Axiomen, \textbf{L$_0$} bis \textbf{L$_8$}, beschreibt, wie wir die Junktoren verwenden wollen. Jedes einzelne Axiom ist eigentlich ein Axiomenschema, weil diese Sätze für beliebige $\mathcal{L}$-Formeln gelten. \textbf{L$_0$} zum Beispiel ist $\varphi \vee \neg\varphi$ --- das Axiom vom ausgeschlossenen Dritten. Es besagt: Es regnet, oder es regnet nicht. \textbf{L$_1$} sagt: Wenn $\varphi$ gilt, dann ist es egal, was $\psi$ ist --- $\varphi$ gilt immer noch.
\end{spoken-clean}

\begin{meta-note}[Projizierter Inhalt: Die logischen Axiome (Teil 1)]
Der Dozent zeigt eine Folie mit den Axiomenschemata $\mathbf{L_0}$ bis $\mathbf{L_8}$.
\end{meta-note}

\begin{math-stroke}[Die logischen Axiome (Junktoren)]
Sei $\mathcal{L}$ eine Signatur und seien $\varphi, \psi, \chi$ beliebige $\mathcal{L}$-Formeln:
\begin{align}
    \mathbf{L_0}\colon &\quad \varphi \vee \neg \varphi && \text{(\textbf{Tertium non datur})} \nonumber \\
    \mathbf{L_1}\colon &\quad \varphi \rightarrow (\psi \rightarrow \varphi) \nonumber \\
    \mathbf{L_2}\colon &\quad (\varphi \rightarrow (\psi \rightarrow \chi)) \rightarrow ((\varphi \rightarrow \psi) \rightarrow (\varphi \rightarrow \chi)) \nonumber \\
    \mathbf{L_3}\colon &\quad (\varphi \wedge \psi) \rightarrow \varphi \nonumber \\
    \mathbf{L_4}\colon &\quad (\varphi \wedge \psi) \rightarrow \psi \nonumber \\
    \mathbf{L_5}\colon &\quad \varphi \rightarrow (\psi \rightarrow (\varphi \wedge \psi)) \nonumber \\
    \mathbf{L_6}\colon &\quad \varphi \rightarrow (\varphi \vee \psi) \nonumber \\
    \mathbf{L_7}\colon &\quad \psi \rightarrow (\varphi \vee \psi) \nonumber \\
    \mathbf{L_8}\colon &\quad (\varphi \rightarrow \chi) \rightarrow ((\psi \rightarrow \chi) \rightarrow ((\varphi \vee \psi) \rightarrow \chi)) \nonumber
\end{align}
\end{math-stroke}

\begin{spoken-clean}[01:25:27 - 01:26:27]
Dann gibt es weitere Axiome, die die Quantoren regeln. \textbf{L$_9$} und \textbf{L$_{10}$} besagen: Wenn etwas für alle gilt, dann gilt es auch für einen speziellen Term --- sofern die Substitution zulässig ist. Und wenn es für einen Term gilt, dann existiert mindestens ein Objekt, für das es gilt.

Zuletzt gibt es noch die Gleichheitsaxiome \textbf{L$_{13}$} bis \textbf{L$_{15}$}. \textbf{L$_{13}$} besagt $\tau = \tau$ --- das ist die Reflexivität. Die anderen regeln die Substitutivität für Relationen und Funktionen.

Das sind unsere Axiome. Ich würde von Ihnen erwarten, dass Sie sich hinsetzen und sich diese anschauen. Ich habe dazu eine Aufgabe auf das Übungsblatt gepackt: Versuchen Sie, diesen Formeln einen Alltagssinn zu geben. Nächste Woche werden wir sehen, wie wir mithilfe dieser Axiome und logischer Schlussfolgerungen Beweise führen können. Vielen Dank fürs Kommen und bis nächste Woche!
\end{spoken-clean}

\begin{math-stroke}[Die logischen Axiome (Quantoren und Gleichheit)]
\textbf{II. Quantorenaxiome (falls Substitution zulässig):}
\begin{align}
    \mathbf{L_9}\colon &\quad \forall \nu \varphi(\nu) \rightarrow \varphi(\tau) \nonumber \\
    \mathbf{L_{10}}\colon &\quad \varphi(\tau) \rightarrow \exists \nu \varphi(\nu) \nonumber \\
    \mathbf{L_{11}}\colon &\quad \forall \nu (\psi \rightarrow \varphi) \rightarrow (\psi \rightarrow \forall \nu \varphi) && (\nu \notin \operatorname{frei}(\psi)) \nonumber \\
    \mathbf{L_{12}}\colon &\quad \forall \nu (\varphi \rightarrow \psi) \rightarrow (\exists \nu \varphi \rightarrow \psi) && (\nu \notin \operatorname{frei}(\psi)) \nonumber
\end{align}

\begin{explanation-of-steps}[Notation $\varphi(\nu)$ und $\varphi(\tau)$]
In den Axiomen $\mathbf{L_9}$ und $\mathbf{L_{10}}$ wird eine informelle, aber sehr gebräuchliche Kurzschreibweise verwendet: $\varphi(\nu)$ deutet an, dass die Variable $\nu$ in $\varphi$ frei vorkommt, und $\varphi(\tau)$ steht für die zulässige Substitution $\varphi(\nu/\tau)$.
\end{explanation-of-steps}

\textbf{III. Gleichheitsaxiome:}
\begin{align}
    \mathbf{L_{13}}\colon &\quad \tau = \tau && \text{(\textbf{Reflexivität})} \nonumber \\
    \mathbf{L_{14}}\colon &\quad \tau_1 = \tau_1' \wedge \dots \wedge \tau_n = \tau_n' \rightarrow (R\tau_1 \dots \tau_n \rightarrow R\tau_1' \dots \tau_n') \nonumber \\
    \mathbf{L_{15}}\colon &\quad \tau_1 = \tau_1' \wedge \dots \wedge \tau_n = \tau_n' \rightarrow (F\tau_1 \dots \tau_n = F\tau_1' \dots \tau_n') \nonumber
\end{align}
\end{math-stroke}

\begin{ai-note}[{Nummerierung der logischen Axiome --- ab Woche 2 um eins verschoben --- Opus 5}]
% [PERSISTENT]
Die obige Liste enthält $\mathbf{L_0}$--$\mathbf{L_{15}}$. Ab \cref{ax:aussagenlogisch} (Woche 2) wird dieselbe Axiomenliste jedoch mit einem zusätzlichen aussagenlogischen Schema geführt,
\[ \mathbf{L_9}\colon \quad \neg \varphi \rightarrow (\varphi \rightarrow \psi) \qquad \text{(\emph{ex falso quodlibet})}, \]
wodurch sich alle nachfolgenden Nummern um eins verschieben: die Quantorenaxiome heissen dort $\mathbf{L_{10}}$--$\mathbf{L_{13}}$ (statt $\mathbf{L_9}$--$\mathbf{L_{12}}$) und die Gleichheitsaxiome $\mathbf{L_{14}}$--$\mathbf{L_{16}}$ (statt $\mathbf{L_{13}}$--$\mathbf{L_{15}}$). Insbesondere ist die Reflexivität $\tau = \tau$ hier $\mathbf{L_{13}}$, ab Woche 2 aber $\mathbf{L_{14}}$. Für alle späteren Kapitel gilt durchgehend die Zählung aus Woche 2.
\end{ai-note}

\begin{nice-box}[Summary of Key Concepts]
In dieser Einführungsvorlesung wurden die administrativen Rahmenbedingungen geklärt und der Grundstein für die formale Logik gelegt. Die zentralen Konzepte waren:
\begin{itemize}
    \item \textbf{Syntax der Prädikatenlogik:} Unterscheidung zwischen logischen Symbolen (Variablen, Junktoren, Quantoren, Gleichheit) und nicht-logischen Symbolen (Konstanten, Funktionen, Relationen), die zusammen die \newterm{Signatur} bilden.
    \item \textbf{Terme und Formeln:} Induktive Definition von \newterm{Termen} (Objektrepräsentationen) und \newterm{Formeln} (Aussagen), wobei \newterm{atomare Formeln} die Basis bilden.
    \item \textbf{Variablenbindung:} Definition von \newterm{freien} und \newterm{gebundenen} Variablen sowie der \newterm{Substitution}, wobei die \newterm{Zulässigkeit} (Vermeidung von Variablen-Einfang) entscheidend ist.
    \item \textbf{Logische Axiome:} Einführung der Hilbert-Axiome (\textbf{L$_0$}--\textbf{L$_{15}$}) als formales Fundament für mathematische Beweise.
\end{itemize}
\end{nice-box}

\begin{ai-global-state-checkpoint-invisible-content}
% timestamp: 01:26:27
% topic: Logische Axiome und Abschluss
% board_state: beamer:axioms-l0-l15
% next_goal: Ende der Vorlesung, Ausblick auf Beweistheorie
% open_loops: none
\end{ai-global-state-checkpoint-invisible-content}

% [SYSTEM] Refinement complete
